\documentclass[12pt]{article}
\usepackage[T1]{fontenc}
\usepackage[utf8]{inputenc}
\usepackage{lmodern}
\usepackage{textcomp}
\usepackage{enumitem}
\usepackage{geometry}
\geometry{margin=1in}
\begin{document}
\begin{center}
Answer Key - History - K-12 Level Assessment 11 \
\small (Answers are ordered exactly as in the question paper.)
\end{center}

\section*{Multiple Choice}
\begin{enumerate}[label=\arabic*.]
\item \textbf{Answer:} B
\\ \textit{Options:} A) north; B) south; C) east; D) west
\\ \textit{Note:} Upper Egypt is south of Lower Egypt because the Nile River flows from south to north, making the southern region historically referred to as Upper Egypt and the northern region as Lower Egypt.
\item \textbf{Answer:} C
\\ \textit{Options:} A) 300 cc.; B) 450 cc.; C) 700 cc.; D) 1200 cc.
\\ \textit{Note:} The mean size of a Homo habilis brain is approximately 700 cc, which reflects its evolutionary development as a key hominin species with a brain capacity larger than that of its predecessors.
\item \textbf{Answer:} A
\\ \textit{Options:} A) colder; Pleistocene epoch.; B) hotter; Miocene epoch.; C) colder; Pliocene epoch.; D) hotter; Pleistocene epoch.
\\ \textit{Note:} The correct answer is based on the fact that during the transition from the Pleistocene epoch, which began around 2.6 million years ago, global temperatures significantly dropped, leading to a series of ice ages characterized by colder climates.
\item \textbf{Answer:} B
\\ \textit{Options:} A) pictograph.; B) petroglyph.; C) pictorial.; D) petrification.
\\ \textit{Note:} A petroglyph is defined as a rock engraving or carving made by removing part of a rock surface, distinguishing it from painted designs.
\item \textbf{Answer:} C
\\ \textit{Options:} A) Mousterian technology produced more projectile points.; B) Mousterian technology produced more usable blade surface.; C) Aurignacian technology produced more usable blade surface.; D) Aurignacian technology was focused on the production of hand axes.
\\ \textit{Note:} Aurignacian technology is characterized by the development of long, thin blades that offered a greater usable surface area for cutting and processing materials, in contrast to the thicker and less efficient tools of Mousterian technology.
\end{enumerate}

\section*{True / False}
\begin{enumerate}[label=\arabic*.]
\setcounter{enumi}{5}
\item \textbf{Answer:} False
\\ \textit{Options:} A) True; B) False
\\ \textit{Note:} Stonehenge is characterized by its stone circle and megalithic structures rather than concentric ridges, and it dates back to around 3000-2000 BCE, which is later than the construction of Poverty Point, approximately 1600-500 BCE.
\item \textbf{Answer:} True
\\ \textit{Options:} A) True; B) False
\\ \textit{Note:} The Bering Strait was a land bridge known as Beringia during the last Ice Age, which connected Asia and North America, facilitating migration between the continents.
\item \textbf{Answer:} True
\\ \textit{Options:} A) True; B) False
\\ \textit{Note:} The Inca political strategy primarily focused on integration and diplomacy rather than relying on ritual sacrifice, as they believed that the sun god could be appeased through offerings and public works rather than through violent means.
\item \textbf{Answer:} True
\\ \textit{Options:} A) True; B) False
\\ \textit{Note:} The Pinnacle Point site in South Africa has been archaeologically dated to around 71,000 years ago, revealing a sophisticated sequence of microblade production and usage by early humans, which demonstrates advanced tool-making skills and social organization.
\item \textbf{Answer:} True
\\ \textit{Options:} A) True; B) False
\\ \textit{Note:} The statement is true because fossil and genetic evidence indicate that anatomically modern Homo sapiens first emerged in Africa approximately 200,000 years ago and have survived to the present day.
\end{enumerate}

\section*{Long Form Response}
\begin{enumerate}[label=\arabic*.]
\setcounter{enumi}{10}
\item \textbf{Answer:} Bipedal locomotion, or walking on two legs, is significant in the evolution of hominids because it allowed early humans to travel long distances efficiently and free their hands for tool use and carrying objects. This adaptation provided several advantages, such as improved visibility in open landscapes and the ability to gather food and resources more effectively. Compared to other species that remained quadrupedal, bipedal hominids could navigate diverse environments, leading to greater survival and adaptability. Overall, bipedalism played a crucial role in shaping the physical and social development of early human ancestors, setting them apart from other primates and contributing to the rise of modern humans.
\\ \textit{Note:} correct because it highlights how bipedal locomotion enabled early hominids to travel efficiently, use tools, and adapt to various environments, providing them with survival advantages over quadrupedal species.
\item \textbf{Answer:} Ancient Egypt had a societal structure centered around the pharaoh, who was considered both a political and religious leader. The population of ancient Egypt is estimated to have been around 3,000,000 people, with approximately 75\% of them being farmers. These farmers were crucial to the economy and daily life, as they worked the fertile land along the Nile River, growing crops that sustained the population and supported the pharaoh's projects. Their labor ensured food security and contributed to the wealth of the state, highlighting the vital role of agriculture in maintaining the stability and prosperity of ancient Egyptian society.
\\ \textit{Note:} The answer correctly identifies the central role of the pharaoh in ancient Egyptian society and emphasizes the significant proportion of the population that were farmers, highlighting their essential contributions to the economy and sustenance of the civilization.
\end{enumerate}

\vfill
\noindent\textit{End of Answer Key}
\end{document}